\section*{AI-use statement}

This paper was drafted by a crew of AI research-assistant agents
(Free-Research-Crew, protocol v0.2) under the direction of the human author, Bryce Wang.
He approved the research question (2026-11-20, 16:53 UTC), the design (2026-11-20,
17:49 UTC) and the results (2026-11-21, 02:17 UTC), each in chat after
reading the Reviewer's verdict. The manuscript itself is subject to the submission gate. He takes responsibility for the paper's content.

% Counts are macros from process.tex (`crew.py numbers`); who did what is taken from
% log.jsonl and reviews/.
\begin{itemize}
  \item \textbf{Roles and models.} Every role (PI, Scout, Designer, Data RA, Econometrician,
        Figures RA, Writer, Reviewer and Replicator) runs on Claude Opus 5.5
        (\texttt{claude-opus-5-5}) in Claude Code, each role as a separate agent. The PI
        drafted the question and hypotheses, recorded deviations D3--D8, and decided D7 and, with the human author's consent, D8; the Scout
        found and verified the sources and transcribed the published targets; the Designer
        wrote the design; the Data RA fetched and prepared the data and, at the paper stage, wrote and
        executed the descriptive run behind the summary-statistics table (D7); the Econometrician
        wrote and tested the estimation code, executed the analysis runs, including the
        post hoc diagnostic and the SAS-listing check, and wrote the result claims; the
        Figures RA drew the figure; the Writer drafted this paper from the claims; the
        Replicator ran the clean rerun. The PI's brief to each role is
        archived in \texttt{briefs/}.
  \item \textbf{Review.} The Reviewer ran \CrewReviewRounds{} rounds on
        \CrewReviewerModels{} (\CrewReviewIndependence{}: a new context in every round, the
        same model family as the other roles): \CrewRoundsQuestion{} on the question,
        \CrewRoundsLiterature{} on the literature, \CrewRoundsDesign{} on the design,
        \CrewRoundsCode{} on the code, \CrewRoundsResults{} on the results and
        \CrewRoundsPaper{} on the paper. It raised \CrewFindingsCritical{} critical,
        \CrewFindingsMajor{} major and \CrewFindingsMinor{} minor distinct findings (a
        finding carried into a later round is counted once, at its latest severity). The last round of each completed stage passed, with minor findings answered
        before the next gate. A code review passed before every analysis run except the D8 figure rerun, including
        separate rounds for the post hoc diagnostic, the SAS-listing check and the D7
        descriptive run.
  \item \textbf{Runs.} \CrewRuns{} runs in the ledger (\CrewRunsAnalysis{} analysis,
        \CrewRunsDataPrep{} data preparation); \CrewRunsSuperseded{} were replaced by a
        later run with the same id and archived (data-preparation and figure runs only; no run
        behind a result claim was replaced); \CrewRunsUnreviewed{} analysis
        run was not covered by a passing code review: the figure run repeated at the paper
        stage to correct its axis-label units only (deviation D8), with the human author's
        consent in chat, after the four code-review rounds had been used.
  \item \textbf{Replication.} The Replicator's clean rerun (\texttt{crew.py rerun}, in a
        fresh environment built from \texttt{code/requirements.txt}) has status
        \CrewRerunStatus{}: \CrewRerunReproduced{} of \CrewRerunTotal{} runs reproduced.
        Whether the work can be rerun from scratch is what this record states, not
        something the authoring roles assert.
  \item \textbf{Findings accepted over the Reviewer's objection.} None.
  \item \textbf{Evidence.} Every result number in the text, the tables and the figure caption
        is a macro generated from the metrics of a logged run. The numbers typed by hand are
        design facts fixed before any analysis run (critical values, targets, the
        pre-registered minimum detectable difference), published numbers quoted next to
        their replication, numbers from the literature quoted as the Scout's source notes
        record them (for example first-stage $F$ statistics, intervals and direct-effect
        estimates of earlier studies), and process facts such as dates. Card's SAS listing,
        which holds regression output, was left unopened until every claim of the approved
        plan was written. Claims added later came from metrics already written by earlier
        runs (results review), from the D7 descriptive run, or, at the paper review, from
        metrics already written (target and item counts, including the listing run's own,
        the KWW sample size, and the post hoc, descriptive homoskedastic statistic at zero);
        some claim texts were reworded at the paper stage. No claim value changed after the
        listing was read. Every
        citation was checked against its source by the Scout; the verification method for
        each entry (full text, abstract or DOI resolution) is recorded in
        \texttt{sources/sources.json}, and the paper cites each source for no more than was
        read. The full workspace, with the plan, run ledger and notes, code, data manifest,
        reviews, briefs and source notes, is at
        \url{https://github.com/brycewang-stanford/Free-Research-Crew/tree/main/gallery/ss-card-proximity}.
\end{itemize}
